% Template for post or presentation abstracts submitted to ILAS 2022
% 1. Fill in the presenter's information (name, affiliation, email
% address) and talk information (title of presentation, abstract).
% 2. Compile to make sure there are no errors.
% 3. Save the .tex file for your abstract with a distinctive name,
% ideally "FamilynameGivenname.tex".
% 4. Upload your .tex file to https://clr.ie/129557
%       
% * Please keep your LaTeX commands simple, and avoid the use of
%    user-defined macros.
% * Include details of co-authors and funding acknowledgements after the abstract.
% * Upload only the LaTeX file (with .tex extension).
% * Questions: Contact Niall Madden (Niall.Madden@NUIGalway.ie)

\documentclass[a4paper]{amsart} 
\usepackage{amssymb} 
\usepackage{hyperref}
\pagestyle{empty}
\date{} 

%% IGNORE THE NEXT 4 LINES
\makeatletter % 
\def\labelenumi{\theenumi.}
\def\theenumi{\@arabic\c@enumi}
\makeatother
%% STOP IGNORING NOW

\begin{document}

\title{A scalable and robust vertex-star relaxation for high-order FEM}
\author{Patrick E.~Farrell} %Presenting author only
\address{University of Oxford} % Affiliation of presenting author
\email{patrick.farrell@maths.ox.ac.uk} % Email address of presenting author only

\maketitle

High-order finite element methods (FEM) offer numerous advantages. They are especially attractive on modern supercomputers, due to their arithmetic intensity and rapid convergence for smooth solutions. However, all aspects of a code must change at high-order, from the choice of basis functions to matrix-free assembly strategies and on to postprocessing and visualisation. A particularly important challenge is to develop preconditioners that operate matrix-free, without ever accessing even the local tensor for a single cell.

One promising strategy for preconditioners for high-order FEM is the use of $p$-multigrid.
Pavarino proved that the two-level method with vertex patch relaxation for the high-degree problem and a
low-order coarse space gives a solver that is robust in polynomial degree for symmetric and coercive
problems \cite{pavarino93_farrell}.
However, for very high polynomial degree it is
not feasible to assemble or factorize the matrices for each vertex patch, since they are dense. 

In this work we introduce a direct solver for separable vertex patch problems that scales
to very high polynomial degree on tensor product cells. The solver
constructs a carefully-chosen tensor product basis that diagonalizes the blocks in the
stiffness matrix for the internal degrees of freedom of each individual
cell. As a result, the non-zero structure of the cell matrices is that of
the graph connecting internal degrees of freedom to their projection onto
the facets. In the new basis, the patch problem is as sparse as a low-order
finite difference discretization, while having a sparser Cholesky
factorization.  We can thus afford to assemble and factorize the matrices
for the vertex-patch problems, even for very high polynomial degree. In turn, this enables the use of
fast $p$-multigrid solvers. In the
non-separable case, the method can be applied as a preconditioner by
approximating the problem with a separable surrogate.

We demonstrate the approach by solving the Poisson equation 
and a $H(\mathrm{div})$-conforming interior penalty discretization of 
linear elasticity in two dimensions at polynomial degree $p=31$ and in three dimensions at $p = 15$.

%% Coauthor details here (optional - delete if not applicable)
% and funding acknowledgment (optional - delete if not applicable)
\emph{This is joint work with Pablo D.~Brubeck (Oxford).
Supported by a Mathematical Institute departmental scholarship, and the Engineering and Physical Sciences Research Council, grants EP/R029423/1 and EP/W026163/1.}

\begin{thebibliography}{1}
\bibitem{pavarino93_farrell}
L.~F.~Pavarino.
\newblock Additive {S}chwarz methods for the $p$-version finite element method.
\newblock {\em Numer. Math.} 66:493–515, (1993).
\end{thebibliography}


\end{document}


% Please leave the next bit alone  
%%% Local Variables: 
%%% mode: latex
%%% TeX-master: "../ILAS-2022"
%%% End:
